\section{}

\comment{
Alternative clustering approaches
I think the authors should include the discussion of alternative clustering methods
in the main text (or in the appendix).
}

\reply{
  We thank the reviewer for the suggestion. In the revised manuscript, we have added a discussion of alternative clustering approaches. We also clarified why the $k$-medoids algorithm was selected for the main analysis. Specifically, the $k$-medoids algorithm has the advantage that each cluster center is an actual player rather than an artificial average point. This makes the resulting clusters easier to interpret in this context, because each cluster can be summarized by a representative player. 

  We have also clarified that the choice of the method is not intrinsic to the proposed framework. The representation provides a general low-dimensional summary of each player's shooting pattern and different clustering algorithms (or other statistical methods) could be applied to the representation depending on the goal of the analysis. For example, hierarchical clustering is useful when we want to study nested similarities between players; mixture models are useful when a probabilistic assignment is desired and a density-based method is useful to identify outliers.

  We have added the following sentences to the Discussion section:

\begin{quote}
\textcolor{blue}{Although, we use the $k$-medoids clustering algorithm in our analysis, the MFPCA representation is not tied to a particular clustering algorithm. Actually, once each player's shot chart has been summarized by a finite vector of scores, all of the standard clustering algorithms could be applied. For example, $k$-means clustering could be used to identify groups around average score profiles, hierarchical clustering could show nested similarities among players, mixture models could provide probabilistic cluster memberships and density-based methods could be use in case of outliers. We use the $k$-medoids algorithm because the center of each cluster is an actual player from the NBA, rather than an artificial mean profile. This provides a direct basketball interpretation of each cluster through its medoid.}
\end{quote}
}

\comment{
Since the data span multiple seasons (2018-19 through 2023-24), shot selection can
change across seasons (and also within a season), due to numerous reasons (e.g.,
changes in team, coaching, usage rate, playing style, etc.). This is perhaps outside
the scope of the paper, but I’m interested in seeing a discussion on whether your
method can be extended to handle this. I think having a couple of paragraphs on
this in the Discussion would make the paper more intriguing to interested readers.
For example, young players may enter the league without a reliable 3-point shot and
therefore prefer scoring near the basket. Over time, they may develop a consistent
3-point shot and become more comfortable taking 3s. A similar pattern may apply
to centers. Players such as Al Horford and Brook Lopez once favored post-ups and
shots around the rim but later expanded their range and became good at 3-point
shooting. (These are just examples of the top of my head; these developments by
Horford and Lopez occurred before the sample of seasons considered in this paper.)
Alternatively, have you also considered clustering at the player-season level instead
of just one single aggregated shot chart for each player? Extending this to a
player-season level could further validate your finding of empirical clusters having
low agreement with official positions. For instance, it could help quantify the
role/positional evolution of a player (related to what I described above).
}

\reply{

  We thank the reviewer for this suggestion. We agree that aggregating shots across multiple seasons, while providing a stable summary, may hide potentially meaningful temporal changes. A similar procedure could be applied after redefining each observational unit from ``player'' to ``player-season''. The resulting MFPCA scores could then be analyzed longitudinally, considering that the observations are not independant. We could then model score trajectories as a function of age or something else. Clustering the player-season scores would also allow players to move between shooting-style groups over time and quantifying role evolution. While being outside the scope of the paper, we added a paragraph in the Discussion section explaining how the proposed method could be extended to account for season-to-season variability.

We have added the following sentences to the Discussion section:

  \begin{quote}
\textcolor{blue}{We aggregated shooting behavior across the 2018-2019 through 2023-2024 regular season. However, this aggregation may mask meaningful season-to-season changes in shooting behavior, as a player's shot selection can evolve with time due to development, injuries, coaching strategy, etc. A natural extension of our framework would be to perform the analysis at the player-season level rather than at the aggregated player level. In this setting, the MFPCA scores could be used to represents the evolution of shooting-style over time. This could help to quantify the temporal evolution of a player's role, identify players whose shooting patterns change across seasons, and examine whether these changes correspond to team changes, coaching changes or something else. This could also strengthen the comparison between empirical shooting-pattern clusters and conventional NBA positions. We could assign a cluster to each player-season and study the stability of these assignments.}
  \end{quote}
  {\color{purple}Great reply but just TBC on last point\dots think we have it wrong way around (allocate player-season to cluster rather than cluster to player-seasondots)}
}
